% !TeX root = ../papers/nc-p3.tex
\section{有限な階数一切断のArtin環と三方向の分離}
\label{sec:finite-marked-artin-secants}
\begin{theorem}\label{thm:finite-marked-artin-secants}
共通切断 $E$ の $\Sigma$ が長さ3のGorensteinスキームなら、
接次元0または1の印でその側の反対対称性が成り立つ。三重の一点も含む。
長さ4なら $B=H^0(\mathcal O_\Sigma)$ は必ずGorensteinであり、
非退化な積跡 $\tau:B\to k$ を選ぶと、左右の三次元基底変更で $E$ は
\[
       z\longmapsto \big((a,b)\longmapsto\tau(zab)\big),
       \qquad z\in B,\quad a,b\in\ker\tau
\]
に一致する。特に $E$ のこの行列空間としての同型型は $B$ で決まる。
この部分切断の対称表示だけでは四窓の対称性はまだ従わない。
\end{theorem}
\begin{proof}[射影と積の表示]
$\Sigma$ のどの長さ3の部分も、どちらの射影でも一本の直線に入らない。
仮に入れば $E$ の非零片側縮約の核に全体が入り、この核は次元2以下である。
その射影直線の全二次小行列式は長さ3で消えるので直線全体で消え、有限性に反する。
直線束を自明化して二つの三座標の像を $U,V\subset B$ とする。
長さ3なら $U=V=B$ で、$\langle\Sigma\rangle$ のベクトル空間は
$B\otimes B\to B$ の積の双対像となる。
長さ4なら $U,V$ は超平面であり、どちらも非零イデアルを含まない。
実際イデアルは長さ1のイデアルを含み、対応する長さ3の商に座標空間の三次元性が必要である。
よって $U,V$ を零化する二つの汎関数はともに完全積跡であり、その比は単元倍である。
一方の自明化を変えて $U=V=\ker\tau$ とする。積 $UV$ は $B$ を張る。
さもなくば $0\ne z\in B$ について $zU\subset U^\perp=k1$ となる。
単元なら次元に反し、非単元なら $zB$ は真のイデアルなので $zU=0$。
このとき $U$ は $\operatorname{ann}(z)$ に含まれるが、$\operatorname{ann}(z)\ne0$ の場合
$U+\operatorname{ann}(z)=B$（$U$ は非零イデアルを含まない）なので矛盾する。
従って $\langle\Sigma\rangle=\mathbb P(E)$、積の双対で主張の表示を得る。
他の完全跡は単元を掛けた跡であり、二つの引数と $z$ に単元を配分すれば同型になる。
逆向きにも行列空間から階数1切断の環が復元される。

\end{proof}
\begin{proof}[長さ3の相対割線]
$B$ の完全跡を固定する。局所Artin基底 $S$ 上で $\Sigma$ の次数2の部分を指定することは、
$B_S=B\otimes S$ に直和因子の直線イデアル $I$ を指定することと同じである。
$B$ は $I$ に指標 $\chi:B\to S$ で作用する。$I$ は $\ker(\chi_S)$ の
完全跡による直交補、すなわちこの核の零化イデアルに一致する。
実際そこに含まれ、どちらも直和因子の直線である。逆に任意の指標がこの直線を与える。
相対割線に対応する二次元行列空間は $I$ の零化汎関数、すなわち完全跡で同定すれば
$\ker(\chi_S)$ の元による乗法である。その共通左核は $I$ である。
従って全ての相対左方向は、$B^*$ 内の指標を完全跡で $B$ に移した
一定の長さ3のスキーム $D_3\subset\mathbb P^2$ を通る。
これは $B$ の完全一次系での埋込みであり、その三座標は自明化後に $B$ そのものを張る。

定理\ref{thm:finite-marked-secant-pencil}の差の局所消去をもう一度用いる。
印以外で差が非零なら $Z=f_0 V$ とし、$\operatorname{ann}(f_0)$ を法として相対割線を得る。
その行列式二次で切る平坦族は上に述べた次数2の部分である。
よって位置1の印から射影した三座標 $l$ に対して、$D_3$ で消える二次式 $Q$ は $Q(l)Z=0$ を満たす。
曲線状の印では $Z=t^sV$ のとき、$l/t$ の方向が $t^{m-s-1}$ を法として $D_3$ を通る。
従って $Q(l)=t^2 Q(l/t)$ は $t^{m-s}$ を法として零である。
$m-s\le1$ でもこの結論は成り立つ。被約の印で差は零である。
ゆえに各第二辺関係の評価 $f\in C=U_0U_1$ に対し $Q(l)f=0$ が全域で成立する。

このような $f$ は零である。$B$ に長さが一つずつ増えるイデアルの列を取り、
各イデアルを $l$ の三つの一次式の中に座標制限の同型で持ち上げる。
前のイデアルの全ての持上げが $f$ を消すとし、次のイデアルの元を $h$ と書く。
長さ1の商の支持方向で消える一次式 $k$ に対し、$kh$ の値は前のイデアルに入る。
自明化が単位 $1\in B$ を与える一次式を $h_0$ とし、その値と同じ前段の持上げを $h'$ とすれば
$kh-h_0 h'$ は $D_3$ 上で消える二次式である。従って $khf=0$。
これは $hf$ の支持が位置1でその方向と印を結ぶ直線に入ることを意味する。
第一座標のみでこの支持環へ全射することは三つの被約割線の場合と同じである。
$h\in U_1$ と $\tau_X(fhU_0)=0$ により $hf=0$ を得る（$\tau_X$ は長さ20の辺の完全跡）。
帰納して全 $l$ が $f$ を消す。補題\ref{lem:finite-adjacent-word-support}から $f=0$ である。
\end{proof}
